\documentclass[11pt]{article}
\usepackage[utf8]{inputenc}
\usepackage[T1]{fontenc}
\usepackage{amsmath}
\usepackage{amssymb}
\usepackage{amsthm}
\usepackage{hyperref}
\usepackage{listings}

\lstset{
  basicstyle=\ttfamily\small,
  breaklines=true,
  frame=single,
  columns=fullflexible,
  escapechar=@
}

\newtheorem{theorem}{Theorem}
\newtheorem{lemma}{Lemma}
\theoremstyle{definition}
\newtheorem{definition}{Definition}

\title{Disproof of Conjecture 00000007468\\(weight-limsup shift spectral disk)}
\author{Zhiyuan Zhang (GitHub: LibertychaserUS)}
\date{Verdict: FALSE (counterexample verified; Lean 4 formal disproof compiled)}

\begin{document}
\maketitle

\section{The conjecture as filed}

The statement in the competition file \texttt{conjectures/00000007468.md} reads, in English:

\begin{quote}
Definition: The spectra of weighted shifts. Conjecture: The spectral disk of the shift holds,
and the radius of the disk is the limit superior of the weights.
(weight-limsup shift spectral disk)
\end{quote}

\section{Precise reading}

Let $H=\ell^2(\mathbb{N})$ with $\mathbb{N}=\{0,1,2,\dots\}$ and orthonormal basis $(e_k)_{k\ge 0}$.
For a bounded sequence $w=(w_k)_{k\ge 0}$ of complex numbers, the unilateral forward weighted shift
$W=W_w$ is the bounded operator determined by $W e_k = w_k e_{k+1}$, that is,
$(Wx)_0=0$ and $(Wx)_{n+1}=w_n x_n$.

We use the following standard definitions.
\begin{itemize}
\item Spectrum: $\sigma(T)=\{\lambda\in\mathbb{C}:\ \lambda I-T \text{ is not invertible in } B(H)\}$.
\item Spectral radius: $r(T)=\sup\{|\lambda|:\lambda\in\sigma(T)\}$. Gelfand's formula gives $r(T)=\lim_n\|T^n\|^{1/n}$.
\item Upper limit of the weights: $\limsup_k w_k=\lim_{K\to\infty}\sup_{k\ge K}|w_k|$.
\item Spectral disk: a set of the form $\overline{D}(c,r)=\{z:|z-c|\le r\}$.
\end{itemize}

The conjecture is read universally: for every bounded weight sequence $w$ there is a centre $c\in\mathbb{C}$
such that $\sigma(W_w)=\overline{D}(c,\limsup_k|w_k|)$.
The centre is left free, so this is the weakest reading of the disk clause.
Any stronger reading (for example with centre $0$) is refuted a fortiori by the same counterexample.

\begin{definition}
A bounded weight sequence $w$ is a \emph{counterexample} if no $c\in\mathbb{C}$ satisfies
$\sigma(W_w)=\overline{D}(c,\limsup_k|w_k|)$.
\end{definition}

\section{The counterexample and its verification}

Define the weights
\[
w_k=\begin{cases} 2, & k \text{ odd},\\ \tfrac12, & k \text{ even}.\end{cases}
\]
All weights satisfy $|w_k|\le 2$, so $W=W_w$ is bounded with $\|W\|\le 2$.

\begin{lemma}[limsup]
$\limsup_k w_k=2$.
\end{lemma}
\begin{proof}
Every $w_k\le 2$, and $w_k=2$ for every odd $k$, so $\sup_{k\ge K}w_k=2$ for every $K$.
\end{proof}

\begin{lemma}[consecutive products]
$w_k\,w_{k+1}=1$ for every $k\ge 0$. Consequently $W^2e_k=e_{k+2}$, and $W^2$ is an isometry
(unitarily equivalent to $S\oplus S$, where $S$ is the unweighted unilateral shift).
\end{lemma}
\begin{proof}
If $k$ is odd then $w_k=2$ and $w_{k+1}=\tfrac12$; if $k$ is even then $w_k=\tfrac12$ and $w_{k+1}=2$.
In both cases the product is $1$. Then $W^2 e_k=w_k w_{k+1}e_{k+2}=e_{k+2}$.
\end{proof}

\begin{lemma}[powers]
For $n\ge 1$, $\|W^n\|=\sup_{k\ge 0}\prod_{i=k}^{k+n-1}w_i$, which equals $2$ for odd $n$ and $1$ for even $n$.
Hence $\|W^n\|^{1/n}\to 1$ and $r(W)=1$.
\end{lemma}
\begin{proof}
Since $W^ne_k=\big(\prod_{i=k}^{k+n-1}w_i\big)e_{k+n}$ and the vectors $e_{k+n}$ are orthonormal,
$\|W^n\|=\sup_k\prod_{i=k}^{k+n-1}w_i$. For even $n$, the window splits into $n/2$ consecutive pairs
$(w_j w_{j+1})$, each equal to $1$. For odd $n$, the window equals $w_k$ times $(n-1)/2$ pairs
$(w_{j}w_{j+1})$ starting at $j=k+1$, so the product is $w_k$ and the supremum over $k$ is $2$.
Gelfand's formula then gives $r(W)=\lim_n\|W^n\|^{1/n}=1$.
\end{proof}

\begin{lemma}[spectrum lies in the unit disk]
$\sigma(W)\subseteq\overline{D}(0,1)$.
\end{lemma}
\begin{proof}
$\|W^2\|=1$, so $\sigma(W^2)\subseteq\overline{D}(0,\|W^2\|)=\overline{D}(0,1)$.
By the spectral mapping theorem for polynomials, $\sigma(W^2)=\{\lambda^2:\lambda\in\sigma(W)\}$.
Hence $|\lambda|^2\le 1$ for every $\lambda\in\sigma(W)$.
\end{proof}

In fact $\sigma(W)$ equals the closed unit disk. Since $W^2$ is unitarily equivalent to $S\oplus S$ and
$\sigma(S)$ is the closed unit disk, $\sigma(W)^2=\overline{D}(0,1)$, and rotation invariance
($U_t e_k=e^{ikt}e_k$ gives $U_tWU_t^*=e^{it}W$) forces $\sigma(W)=\overline{D}(0,1)$.
This equality is established by the analytic argument printed in \texttt{check.py}; the Lean file
only uses the inclusion $\sigma(W)\subseteq\overline{D}(0,1)$, which suffices.

\begin{theorem}
For $w_k$ as above, $\sigma(W_w)$ is not a closed disk $\overline{D}(c,2)$ for any $c\in\mathbb{C}$.
Hence $w$ is a counterexample to the conjecture.
\end{theorem}
\begin{proof}
Suppose $\sigma(W)=\overline{D}(c,2)$. Then $c+2$ and $c-2$ lie in $\sigma(W)\subseteq\overline{D}(0,1)$,
so $|c+2|\le 1$ and $|c-2|\le 1$. By the triangle inequality
$4=|(c+2)-(c-2)|\le |c+2|+|c-2|\le 2$, a contradiction.
\end{proof}

Numerical checks (exact rational arithmetic, \texttt{check.py}): all $w_kw_{k+1}=1$ for $k<1000$;
$\sup_{k<1000}w_k=2$ and all tail suprema are $2$; for $n=1,\dots,60$ the exact value of $\|W^n\|$ equals
$2$ for odd $n$ and $1$ for even $n$, and the window maximum over $k\in[0,2)$ equals the maximum over $[0,200)$.

\section{Lean formalization}

The file \texttt{Disproof.lean} (Lean 4.33.1, Mathlib v4.33.1, built with \texttt{lake env lean}) contains:
\begin{itemize}
\item \textbf{Hilbert space.} \texttt{Hs := lp (fun \_ : Nat => Complex) 2}, that is, $\ell^2(\mathbb{N})$.
\item \textbf{Shift.} \texttt{shiftCoord w x} implements $(Wx)_0=0$, $(Wx)_{n+1}=w_nx_n$;
\texttt{shiftOp w M hM} (of type $\texttt{Hs}\to_L\texttt{Hs}$, i.e.\ a continuous linear map over $\mathbb{C}$)
is the bounded operator built with \texttt{LinearMap.mkContinuous} (bound $M$, where
\texttt{hM} is the hypothesis $|w_k|\le M$ for all $k$).
\item \textbf{Spectrum.} Mathlib's \texttt{spectrum Complex (shiftOp w M hM)}, the spectrum in the Banach algebra
of bounded operators, written $\sigma(W)$ here.
\item \textbf{Conjecture.} \texttt{ConjectureHolds : Prop} is
\[
\forall\, w\colon\mathbb{N}\to\mathbb{C},\ \forall\, M\in\mathbb{R},\ (\forall k,\ |w_k|\le M)\ \Rightarrow\
\exists\, c\in\mathbb{C},\quad \sigma(W_w)=\overline{D}\Big(c,\ \limsup_k |w_k|\Big).
\]
\item \textbf{Main theorem.} \texttt{TLMC7468.conjecture\_00000007468\_fails}, stated as the negation
$\neg\,\texttt{ConjectureHolds}$.
\end{itemize}

Proof structure in Lean.
\begin{enumerate}
\item \texttt{norm\_sq\_eq\_tsum}: $\|x\|^2=\sum_n|x_n|^2$ (from \texttt{lp.norm\_rpow\_eq\_tsum}).
\item \texttt{shiftOp} is bounded by $M$ (\texttt{norm\_shiftFn\_le}), via reindexing of the tsum.
\item \texttt{tsum\_shift\_shift}: $\sum_n|(W^2x)_n|^2=\sum_n|x_n|^2$, using \texttt{wgt\_succ\_mul}
($w_{n+1}w_n=1$) and \texttt{Summable.tsum\_eq\_zero\_add} twice. Hence $\|W^2\|\le 1$ (\texttt{norm\_WW\_op\_le}).
\item \texttt{spectrum\_shiftOp\_subset}: spectral mapping \texttt{spectrum.map\_pow\_of\_pos} with
\texttt{spectrum.subset\_closedBall\_norm} gives $\sigma(W)\subseteq\overline{D}(0,1)$.
\item \texttt{two\_le\_limsup\_wgt}: $\limsup_k|w_k|\ge 2$, from the frequently-odd argument
and the bound $|w_k|\le 2$.
\item Main theorem: if $\sigma(W)=\overline{D}(c,L)$ with $L\ge 2$, then $c+L$ and $c-L$ lie in the unit disk,
so $2L=|(c+L)-(c-L)|\le 2$, contradicting $L\ge 2$.
\end{enumerate}

Axioms (output of \texttt{\#print axioms}):
\begin{lstlisting}
'TLMC7468.conjecture_00000007468_fails' depends on axioms: [propext, Classical.choice, Quot.sound]
\end{lstlisting}

Scope of the formal statement. The formalization treats the conjecture in its universal reading over all
bounded complex weight sequences, with the disk centre free. The Lean proof uses only that
$\sigma(W)$ is contained in the closed unit disk and that $\limsup\|w_k\|\ge 2$; it does not formalize
the equality $\sigma(W)=\overline{D}(0,1)$ or the exact value $\limsup=2$.

\section{Limitations and risks}

\begin{enumerate}
\item The statement is informal. The natural reading (every bounded unilateral weighted shift has spectrum a
disk of radius $\limsup|w_k|$) is refuted here. If the intended class is restricted (for example convergent
weights, or constant weights), the counterexample lies outside it, since $w_k$ does not converge.
The statement names no class, so the universal reading is used.
\item The disk clause is true for this shift (the spectrum is the closed unit disk); the refutation rests on the
radius clause, which is the substantive assertion of the conjecture.
\item The index convention ($\mathbb{N}$ starting at $0$ or $1$) changes no quantity.
\item The Lean file uses standard Mathlib results: Gelfand-type spectral bounds through
\texttt{spectrum.subset\_closedBall\_norm}, spectral mapping for powers, and \texttt{lp} norm identities.
\item A finite combinatorial core (scaled integers, window maxima, consecutive products) is also checked with plain
\texttt{lean} at \texttt{lean/Check.lean}; it is a supplement, not part of the formal disproof.
\end{enumerate}

\section{Reproduction}

\begin{enumerate}
\item Numerical and exact checks: \texttt{python3 check.py} in
\texttt{LibertychaserUS\_submission\_20261027064413/}. Expected: all assertions pass; the verdict block
prints $r(W)=1\neq 2=\limsup w_k$.
\item Lean disproof: \texttt{cd lean} then
\texttt{lake exe cache get \&\& lake build}. Expected: no errors, and the
\texttt{\#print axioms} line above.
\item This report: \texttt{cd LibertychaserUS\_submission\_20261027064413} then \texttt{tectonic report.tex}.
\end{enumerate}

\end{document}
